\documentclass[11pt,a4paper]{article}

\usepackage[utf8]{inputenc}
\usepackage[T1]{fontenc}
\usepackage{lmodern}
\usepackage[a4paper,margin=2.3cm]{geometry}
\usepackage{amsmath,amssymb,amsthm}
\usepackage{bm}
\usepackage{booktabs}
\usepackage{siunitx}
\usepackage{graphicx}
\usepackage{hyperref}
\usepackage{xcolor}
\usepackage{listings}
\usepackage{microtype}
\usepackage{caption}

\hypersetup{
  colorlinks=true,
  linkcolor=black,
  citecolor=blue,
  urlcolor=blue
}

\lstdefinestyle{pystyle}{
  language=Python,
  basicstyle=\ttfamily\footnotesize,
  keywordstyle=\color{blue!70!black}\bfseries,
  commentstyle=\color{gray!60!black}\itshape,
  stringstyle=\color{red!60!black},
  numbers=none,
  frame=single,
  framerule=0.4pt,
  rulecolor=\color{gray!40},
  showstringspaces=false,
  breaklines=true,
  tabsize=2
}
\lstset{style=pystyle}

\newcommand{\Ma}{\mathit{Ma}}
\newcommand{\Vol}{\mathit{Vol}}
\newcommand{\Swet}{S_{\mathrm{wet}}}
\newcommand{\Sp}{S_{p}}
\newcommand{\Sb}{S_{b}}
\newcommand{\Lw}{L_{w}}
\newcommand{\Ls}{L_{s}}
\newcommand{\dc}{\delta_{c}}

\title{\bfseries Liu--Liu--Wen--Ding 2019 Variable-Mach Osculating Flowfield Waverider\\[4pt]
\large Implementation and Software Architecture}
\author{Computational implementation notes}
\date{}

\begin{document}
\maketitle

\begin{abstract}
This document describes the Python reproduction of the variable-Mach-number
osculating flowfield waverider method of Liu \textit{et al.}~(2019,
\textit{Acta Astronautica} 162, 160--167). The implementation is a
self-contained package, \texttt{liu2019/}, that generates the 3D geometry
and evaluates inviscid aerodynamic coefficients for the paper's baseline
test case (\mbox{$\Ma_{\mathrm{center}}=6$}, $\Ma_{\mathrm{tip}}=13$,
\mbox{$\beta=13^{\circ}$}) and wraps it in a PyQt5 GUI tab integrated into
the existing waverider-design application. We summarise the governing
equations, detail the function of each module, and catalogue every design
parameter together with its geometric effect.
\end{abstract}

\tableofcontents

% ==============================================================
\section{Overview of the Method}
% ==============================================================

A waverider is a hypersonic lifting body whose leading edges are
constrained to lie on a prescribed shock wave so that the underside is
effectively confined to a single-shock compression field. Classical
methods (Nonweiler, Jones--Moore--Pike, Sobieczky \textit{osculating cones})
assume a \emph{single} design Mach number. The method of
Liu~\textit{et~al.} generalises this by allowing the design Mach number
to vary along the spanwise coordinate~$z$, thereby generating a vehicle
that is simultaneously ``shaped by the shock'' across a range of Mach
numbers.

Five construction steps, following the paper's Section~1.1, are:

\begin{enumerate}
\item A spanwise Mach distribution $\Ma(z)$ is prescribed
      (Eq.~\ref{eq:Ma}).
\item Two exit-plane curves are prescribed in the base plane $x=\Lw$:
      the trailing edge of the freestream (upper) surface, $y(z)$
      (Eq.~\ref{eq:yupper}), and the shock-wave curve $y_{s}(z)$
      (Eq.~\ref{eq:ys}).
\item Their coefficients are determined from six design parameters
      ($\beta, \Lw, W, \Ls, y_{5}, y_{6}$) by Eqs.~\ref{eq:d}--\ref{eq:a}
      and~\ref{eq:A}.
\item For every spanwise station a vertical \emph{osculating plane}
      is erected; a Taylor--Maccoll cone solution is placed in each
      plane at the local Mach number, giving a cone half-angle
      $\dc(z)$. The leading edge and the compression-surface trailing
      edge are obtained by straight-line intersection within the plane.
\item The 3D surfaces are lofted from the array of per-plane records.
\end{enumerate}

% ==============================================================
\section{Mathematical Formulation}
% ==============================================================

\subsection{Spanwise Mach distribution}

The design Mach number varies quadratically with~$z$:
\begin{equation}
\boxed{\;\Ma(z) \;=\; m z^{2} + n, \qquad -\tfrac{W}{2} \le z \le +\tfrac{W}{2}\;}
\label{eq:Ma}
\end{equation}
with coefficients chosen so that $\Ma(0)=\Ma_{\mathrm{center}}$ and
$\Ma(\pm W/2)=\Ma_{\mathrm{tip}}$:
\[
 n = \Ma_{\mathrm{center}}, \qquad
 m = \frac{\Ma_{\mathrm{tip}} - \Ma_{\mathrm{center}}}{(W/2)^{2}} .
\]

\subsection{Exit-plane curves}

The freestream-surface trailing edge is cubic:
\begin{equation}
\boxed{\;y(z) \;=\; a z^{3} + b z^{2} + c z + d \;}
\label{eq:yupper}
\end{equation}
and the shock-wave trace in the base plane is a piecewise quartic
with a flat central portion of half-length $\Ls$:
\begin{equation}
y_{s}(z) \;=\;
\begin{cases}
 A\,(z-\Ls)^{4} & z \ge +\Ls,\\
 0 & -\Ls<z<+\Ls,\\
 A\,(z+\Ls)^{4} & z \le -\Ls.
\end{cases}
\label{eq:ys}
\end{equation}

\subsection{Coefficient determination}

Writing $(y_{5},z_{5})$ for the tip point on the upper trailing edge
and prescribing its slope $\delta_{5}$, and similarly $(y_{6},z_{6})$,
$\delta_{6}$ for the centreline apex, the cubic coefficients follow from
continuity and slope conditions:
\begin{align}
 d &= \Lw \tan\beta , \label{eq:d}\\
 c &= d \tan\delta_{6} , \label{eq:c}\\
 b &= \frac{3\,(y_{5} - d - c z_{5})
           \;-\; z_{5}(d\tan\delta_{5} + z_{5} c)}{z_{5}^{2}} , \label{eq:b}\\
 a &= \frac{y_{5} - d - c z_{5} - b z_{5}^{2}}{z_{5}^{3}} . \label{eq:a}
\end{align}
The quartic shock coefficient is
\begin{equation}
 A \;=\; \frac{y_{5}}{(z_{5}-\Ls)^{4}}. \label{eq:A}
\end{equation}

For the paper's test case ($\delta_{5}=\delta_{6}=0$, $\Lw=\SI{6}{m}$,
$\beta=\ang{13}$, $y_{5}=\SI{0.1608}{m}$, $z_{5}=\SI{1.5}{m}$), the
constants evaluate to
\[
 d = \SI{1.3852}{m}, \quad c = 0, \quad b = -\SI{1.6325}{m^{-1}},
 \quad a = +\SI{0.7256}{m^{-2}}, \quad A = \SI{7.755e-2}{m^{-3}}.
\]

\subsection{Oblique-shock relations}

For a two-dimensional wedge, the deflection angle $\theta$ and wave
angle $\beta$ at freestream Mach number $\Ma$ are related through the
classical $\theta$--$\beta$--$\Ma$ equation:
\begin{equation}
 \tan\theta \;=\; \frac{2}{\tan\beta}\,
 \frac{\Ma^{2}\sin^{2}\beta - 1}
      {\Ma^{2}(\gamma + \cos 2\beta) + 2}. \label{eq:thetabeta}
\end{equation}
Post-shock ratios across a normal component
$M_{1n}=\Ma\sin\beta$ are
\begin{align}
 \tfrac{p_{2}}{p_{1}} &= 1 + \tfrac{2\gamma}{\gamma+1}(M_{1n}^{2}-1),\\
 \tfrac{\rho_{2}}{\rho_{1}} &= \tfrac{(\gamma+1) M_{1n}^{2}}{(\gamma-1) M_{1n}^{2} + 2},\\
 \tfrac{T_{2}}{T_{1}} &= \tfrac{p_{2}/p_{1}}{\rho_{2}/\rho_{1}},\\
 \Ma_{2} &= \frac{1}{\sin(\beta-\theta)}
           \sqrt{\frac{(\gamma-1) M_{1n}^{2} + 2}
                      {2\gamma M_{1n}^{2} - (\gamma-1)}}.
\end{align}
The Mach wave angle $\mu=\arcsin(1/\Ma)$ sets the lower bound of
physically admissible $\beta$, and the detachment angle $\beta_{\det}$
(at which $\theta$ is maximal) sets the upper bound of the weak-shock
branch. The implementation invertedly solves~\eqref{eq:thetabeta} with
Brent's method on $(\mu,\beta_{\det})$.

\subsection{Taylor--Maccoll cone flow}

For axisymmetric conical flow downstream of a conical shock,
the non-dimensional velocity components
$V_{r}(\theta), V_{\theta}(\theta)$ satisfy the Taylor--Maccoll ODE
\begin{equation}
\begin{aligned}
 \frac{dV_{r}}{d\theta} &= V_{\theta}, \\
 \frac{dV_{\theta}}{d\theta} &=
   \frac{V_{\theta} V_{r} V_{\theta}
         - \mathcal{A}\!\left(2 V_{r} + V_{\theta}\cot\theta\right)}
        {\mathcal{A} - V_{\theta}^{2}},\\
 \mathcal{A} &\equiv \tfrac{\gamma-1}{2}(1 - V_{r}^{2} - V_{\theta}^{2}) .
\end{aligned}
\label{eq:TM}
\end{equation}
Given a free-stream $\Ma$ and a desired conical shock angle $\beta$,
Eq.~\eqref{eq:TM} is integrated inward from $\theta=\beta$ with
post-shock initial conditions obtained from the oblique-shock relations.
The integration is terminated when $V_{\theta}=0$; the corresponding
$\theta$ is the cone half-angle $\dc$. The local cone-surface streamline
is a straight line at angle $\dc$ below the freestream direction.

For the paper's test case ($\beta=\ang{13}$) the solver returns
$\dc(\Ma{=}6)=\ang{8.39}$, $\dc(\Ma{=}8)=\ang{10.02}$,
$\dc(\Ma{=}10)=\ang{10.71}$, $\dc(\Ma{=}13)=\ang{11.19}$.

\subsection{Osculating-plane geometry}

At every spanwise station $z$, a vertical $(x,y)$ plane is erected
(the ``osculating plane''). Within it the reduction to 2D gives:
\begin{align}
 (x,y)_{\mathrm{shock\ foot}} &= (\Lw,\; y_{s}(z)), \\
 (x,y)_{\mathrm{upper\ foot}} &= (\Lw,\; y(z)), \\
 x_{\mathrm{LE}} &= \Lw - \frac{y(z) - y_{s}(z)}{\tan\beta}, \label{eq:xLE}\\
 y_{\mathrm{LE}} &= y(z), \label{eq:yLE}\\
 y_{\mathrm{TE,\,lower}} &= y_{\mathrm{LE}} - (\Lw - x_{\mathrm{LE}})\tan\dc(z). \label{eq:yTE}
\end{align}
Equation~\eqref{eq:xLE} places the leading edge at the intersection of
the shock ray (rising upstream at angle $\beta$ from $y_{s}$) with the
freestream-aligned upper surface.
Equation~\eqref{eq:yTE} follows the cone-surface streamline straight
downstream from the leading edge, dropping at angle $\dc(z)$, until it
reaches the base plane $x=\Lw$.

The radius of curvature of the shock curve in its planar form
$y_{s}(z)$ is
\begin{equation}
 R_{\mathrm{osc}}(z) \;=\; \frac{\bigl[1 + y_{s}'(z)^{2}\bigr]^{3/2}}
                                {\bigl|\,y_{s}''(z)\,\bigr|}
                 \;=\; \frac{\bigl[1 + 16 A^{2}(|z|-\Ls)^{6}\bigr]^{3/2}}
                             {12 A (|z|-\Ls)^{2}}, \quad |z|>\Ls.
 \label{eq:Rosc}
\end{equation}
It is reported as a diagnostic; for $|z|<\Ls$ the shock curve is
straight and $R_{\mathrm{osc}}\to\infty$.

\subsection{Geometric metrics}

With $\mathbf{P}_{u}(x,z)$ and $\mathbf{P}_{\ell}(x,z)$ the sampled upper
and lower surface point clouds (each of size $n_{x}\times n_{z}$),
\begin{align}
 \Vol &= \iint_{\Omega} \bigl[y_{u}(x,z) - y_{\ell}(x,z)\bigr]\, dx\,dz, \\
 \Swet &= \sum_{k} A_{k}^{\mathrm{panel}}
          \quad\text{(upper + lower quads)},\\
 \Sp &= \sum_{k} A_{k}^{\mathrm{proj}\,(x,z)}
          \quad\text{(lower surface projected on $y=0$)},\\
 \Sb &= \text{shoelace area of the base-plane polygon stitched from}\\
     &\quad\;y_{u}(z) \text{ above and } y_{\mathrm{TE}}(z) \text{ below},\\
 \eta &= \Vol^{2/3}/\Swet \quad\text{(volumetric efficiency)}.
\end{align}
All surface sums are vectorised by splitting each structured quad into
two triangles and using the cross-product area formula.

\subsection{Aerodynamic model}

For inviscid hypersonic flow, modified-Newtonian impact theory yields
a local pressure coefficient on windward panels
\begin{equation}
 C_{p} \;=\; C_{p,\max}\,\cos^{2}\theta_{\text{loc}},
 \qquad
 C_{p,\max} \;=\; \frac{2}{\gamma \Ma_{\infty}^{2}}
                   \!\left(\frac{p_{0,2}}{p_{\infty}} - 1\right),
\end{equation}
with $C_{p}=0$ on the shadow side. Here $\cos\theta_{\text{loc}}=
-\hat{\mathbf{n}}\cdot\hat{\mathbf{v}}_{\infty}$ is the impact angle
between the outward panel normal and the freestream direction, and
$p_{0,2}/p_{\infty}$ is the Rayleigh--Pitot stagnation-pressure ratio
behind a normal shock. The non-dimensional force on panel $k$ is
\[
 \mathbf{F}_{k}/q_{\infty} \;=\; -\,C_{p,k}\,A_{k}\,\hat{\mathbf{n}}_{k}.
\]
Integrating over all panels and projecting onto freestream and
lift-direction unit vectors yields $C_{L}, C_{D}, L/D, C_{m,z}$ and the
non-dimensional centre-of-pressure location
$X_{\mathrm{cp}}=(M_{z}/F_{y})/L_{\mathrm{ref}}$.

% ==============================================================
\section{Software Architecture}
% ==============================================================

The package is organised around a single-responsibility principle:
each module exposes only what the next module needs.

\begin{center}
\begin{tabular}{@{}p{0.22\textwidth}p{0.7\textwidth}@{}}
\toprule
Module & Responsibility \\
\midrule
\texttt{config.py}        & Paper parameters, reference metrics, tolerances. \\
\texttt{distributions.py} & $\Ma(z)$, $y(z)$, $y_{s}(z)$ and their coefficients. \\
\texttt{shock.py}         & Oblique shock, cone flow, Taylor--Maccoll solver. \\
\texttt{osculating.py}    & Per-plane geometry records (LE, TE, streamline, $\dc$). \\
\texttt{geometry.py}      & 3D assembly, volume / area integrals, STL export. \\
\texttt{aero.py}          & Modified-Newtonian evaluator for $C_{L}, C_{D}, \ldots$ \\
\texttt{validate.py}      & Runs the paper test case and compares to Table~4 / Fig.~12. \\
\texttt{\_\_main\_\_.py}  & CLI: \texttt{python -m liu2019 --validate --stl out.stl}. \\
\bottomrule
\end{tabular}
\end{center}

\subsection{\texttt{config.py}}

Stores four read-only dictionaries: \texttt{PAPER\_PARAMS},
\texttt{PAPER\_TRAJECTORY} (constant-$q$ altitude / Mach points from
Table~2), \texttt{PAPER\_REFERENCE\_GEOMETRY} (Table~4 values), and
\texttt{PAPER\_REFERENCE\_AERO} (Fig.~12 reading), alongside a
\texttt{TOLERANCES} dictionary of fractional acceptance bands.

\subsection{\texttt{distributions.py}}

Five pure functions mirroring Eqs.~\ref{eq:Ma}, \ref{eq:yupper},
\ref{eq:ys}, \ref{eq:d}--\ref{eq:a}, and \ref{eq:A}. Accepts scalar or
\texttt{numpy} array arguments; vectorised for use in both geometry
generation and GUI plotting.

\subsection{\texttt{shock.py}}

Provides:
\begin{itemize}\setlength\itemsep{0pt}
\item \texttt{theta\_from\_beta\_Ma} -- direct evaluation of
      Eq.~\eqref{eq:thetabeta};
\item \texttt{beta\_from\_theta\_Ma} -- Brent inversion on the weak
      branch, with \texttt{DetachedShockError} when $\theta>\theta_{\max}$;
\item \texttt{oblique\_shock\_ratios} -- full post-shock state dictionary;
\item \texttt{mach\_angle}, \texttt{beta\_detachment} -- bounds used in
      the GUI feasibility check;
\item \texttt{taylor\_maccoll\_cone\_angle} -- integrates the T--M ODE
      with \texttt{scipy.integrate.solve\_ivp}, event-terminates at
      $V_{\theta}=0$. Tolerances are tuned
      (\texttt{rtol=1e-5, atol=1e-8, method="RK45"}) so that a single
      solve completes in \SI{\lesssim 1}{ms}; the full waverider with
      200 stations is under \SI{2}{s}. When
      \texttt{waverider\_generator.flowfield.cone\_angle} is importable
      the module delegates to it; otherwise the embedded fallback is
      used.
\end{itemize}

\subsection{\texttt{osculating.py}}

Declares two dataclasses -- \texttt{OsculatingPlaneData} (per-plane
record: $z$, $\Ma$, $\dc$, $R_{\mathrm{osc}}$, $\mathbf{P}_{\mathrm{LE}}$,
$\mathbf{P}_{\mathrm{TE}}$, streamline samples) and
\texttt{OsculatingPlaneSet} (list-like wrapper carrying the derived
coefficient dictionary $\{a,b,c,d,A\}$). The driver
\texttt{build\_all\_osculating\_planes(params, n\_z, n\_x)} sweeps
$z\in[0,+W/2]$ on \texttt{np.linspace}, evaluates
Eqs.~\eqref{eq:xLE}--\eqref{eq:yTE} at every station, and interpolates
$\dc(\Ma)$ from a pre-computed 15-point grid to amortise the
Taylor--Maccoll cost. Symmetry is exploited: only the starboard half is
stored; mirroring is performed in the geometry-assembly stage.

\subsection{\texttt{geometry.py}}

The \texttt{Liu2019Waverider} class owns the sampled upper/lower
surface arrays ($X_{u}, Y_{u}, Z_{u}$; $X_{\ell}, Y_{\ell}, Z_{\ell}$)
of shape $(n_{x}, n_{z})$. All integrals are vectorised:
\begin{itemize}\setlength\itemsep{0pt}
\item \texttt{volume()} -- trapezoidal rule along $x$ at each $z$-column
      (the $x$-axis varies with $z$ because $x_{\mathrm{LE}}$ does),
      then along $z$.
\item \texttt{wetted\_area()}, \texttt{planform\_area()} -- quad
      triangulation with $\tfrac{1}{2}\|\mathbf{v}_{1}\times\mathbf{v}_{2}\|$
      cross-products, summed over the full panel stack.
\item \texttt{base\_area()} -- shoelace formula on the stitched
      base-plane polygon.
\item \texttt{export\_stl(path)} -- ASCII STL of the wetted surfaces
      with per-triangle outward normals.
\end{itemize}
The convenience builder \texttt{build\_liu2019\_waverider(params,
n\_z, n\_x)} wraps the full pipeline.

\subsection{\texttt{aero.py}}

\texttt{Liu2019AeroEvaluator} triangulates the mirrored surfaces once,
caches their areas / centroids / outward normals, and exposes
\texttt{evaluate(Ma, alpha\_deg)} and
\texttt{evaluate\_paper\_trajectory()}. The evaluator is inviscid and
does not require PySAGAS; trends across the paper's four-point Mach
sweep satisfy the \SI{\pm 10}{\percent} aerodynamic tolerance in
\texttt{config.TOLERANCES} for $C_{L}, C_{D}, L/D$ and \SI{\pm 5}{\percent}
for $X_{\mathrm{cp}}$.

\subsection{GUI tab \texttt{liu2019\_waverider\_tab.py}}

Implemented as a \texttt{QWidget} with a horizontal \texttt{QSplitter}:
\begin{itemize}\setlength\itemsep{0pt}
\item left panel -- geometry / Mach / mesh groups in \texttt{QGroupBox},
      with a live feasibility label (shock attachment, $\Ls<W/2$,
      $y_{5}<y_{6}$);
\item right panel -- six sub-tabs (3D view, base plane, $\Ma(z)$,
      $\dc(z)$, aerodynamics, validation table) plus a metric-card
      strip for $\Vol, \Swet, \Sp, \Sb, \eta$ coloured against the
      paper's tolerances;
\item worker threads \texttt{\_GeometryWorker} and \texttt{\_AeroWorker}
      keep the event loop responsive;
\item action row -- ``Export STL\dots'', ``Run aero (Ma 6, 8, 10, 13)'',
      validation badge.
\end{itemize}
Integration into \texttt{waverider\_gui.py} adds one \texttt{try/except}
import block and one \texttt{addTab} call after the Planar tab; no
existing tabs are modified.

% ==============================================================
\section{Design Parameters and Their Geometric Effect}
% ==============================================================

Table~\ref{tab:params} lists every parameter exposed to the user,
its admissible range, its paper value, and its qualitative effect on
the waverider geometry. The right-hand column states which part of the
construction each parameter controls and which metrics it primarily
moves.

\begin{table}[h]
\caption{Design parameters of the Liu 2019 waverider.}
\label{tab:params}
\centering
\small
\renewcommand{\arraystretch}{1.15}
\begin{tabular}{@{}lllp{0.50\textwidth}@{}}
\toprule
Symbol & Role & Paper value & Geometric effect \\
\midrule
$\beta$        & shock angle           & \ang{13}       &
Sets $d=\Lw\tan\beta$ (Eq.~\ref{eq:d}), the base-plane apex height of
the freestream surface. Larger $\beta$ raises $d$, steepens the shock,
grows~$\Vol$ and $\Sb$, and lifts $C_{L}$. Must exceed
$\mu=\arcsin(1/\Ma)$ and fall below $\beta_{\det}$ for every $\Ma$ in
the span. \\

$\Lw$          & vehicle length        & $\SI{6}{m}$    &
Overall chord and moment reference length. Scales every streamwise
dimension linearly (Eqs.~\ref{eq:d}, \ref{eq:xLE}, \ref{eq:yTE}).
Appears in $\Vol\propto\Lw$, $\Swet\propto\Lw$. \\

$W$            & full span             & $\SI{3}{m}$    &
Full lateral extent $[-W/2,+W/2]$. Fixes $z_{5}=W/2$ and enters
the $\Ma(z)$ polynomial coefficient $m=(\Ma_{t}-\Ma_{c})/(W/2)^{2}$
(Eq.~\ref{eq:Ma}). Wider span flattens the $\Ma(z)$ gradient and
stretches the base-plane curves laterally. \\

$\Ls$          & flat-shock half-width & $\SI{0.3}{m}$  &
Width of the central planar-shock region. Within $|z|<\Ls$ the shock
is a straight line in the base plane ($y_{s}\equiv 0$) and the
osculating-plane radius is infinite. Larger $\Ls$ creates a wider
``cone-derived-like'' centre and reduces the outer curvature zone;
requires $\Ls<W/2$. \\

$y_{5}$        & tip TE height         & \SI{0.1608}{m} &
$y$-coordinate of the upper trailing edge at the tip $z_{5}=W/2$.
Controls both the cubic coefficients $a,b$ (Eqs.~\ref{eq:a}, \ref{eq:b})
and the quartic coefficient $A$ (Eq.~\ref{eq:A}): $A=y_{5}/(z_{5}-\Ls)^{4}$,
so shrinking~$y_{5}$ thins the shock curve at the tips. Must be smaller
than~$y_{6}$. \\

$y_{6}$        & centreline apex height& $\SI{1.608}{m}$&
Reference value for the upper-surface nose height. Does not enter
Eqs.~\ref{eq:d}--\ref{eq:a} (the centreline value of $y(z)$ is fixed
by~$d$) but is kept as a validation anchor and as an upper bound for
$y_{5}$ (constraint $y_{5}<y_{6}$). \\

$\delta_{5}$   & tip TE slope          & $0$            &
Slope $dy/dz$ at the tip. The paper uses zero. Non-zero values appear
in Eqs.~\ref{eq:b} and change the cubic curvature of the upper TE. \\

$\delta_{6}$   & centreline TE slope   & $0$            &
Slope at the centreline. Zero in the paper; the GUI hides this
parameter. Non-zero values set $c=d\tan\delta_{6}\ne 0$ in
Eq.~\ref{eq:c}. \\

$\Ma_{c}$      & centreline $\Ma$      & $6$            &
Governs the minimum design Mach number. Small $\Ma_{c}$ gives a more
slender centre-line cone (small $\dc$ at $z=0$), making the centreline
compression surface shallower. \\

$\Ma_{t}$      & tip $\Ma$             & $13$           &
Governs the maximum design Mach number. Large $\Ma_{t}$ gives wider
cone angles near the tips, deepens the outboard compression surface
and increases~$\Vol$. \\

$\gamma$       & specific-heat ratio   & $1.4$          &
Appears in every gas-dynamic relation (Eqs.~\ref{eq:thetabeta},
\ref{eq:TM}). Not user-exposed in the GUI. \\

$n_{z}$        & spanwise planes       & $200$          &
Number of osculating-plane stations on the half-span. Sets the
resolution of the leading edge and base contour. Larger values
converge every integrated metric roughly as $\mathcal{O}(n_{z}^{-2})$. \\

$n_{x}$        & streamwise strips     & $100$          &
Number of samples along each cone-surface streamline. Sets the
resolution of $\Swet$ and $\Vol$ integrals; convergence is
$\mathcal{O}(n_{x}^{-2})$. \\
\bottomrule
\end{tabular}
\end{table}

\paragraph{Qualitative sensitivity map.}
\begin{itemize}
\item Increasing $\beta$ and/or $\Ma_{t}$ deepens the compression
      surface $\Rightarrow$ $\Vol{\uparrow}$, $\Swet{\uparrow}$,
      $C_{L}{\uparrow}$, $C_{D}{\uparrow}$, $L/D$ typically decreases.
\item Increasing $\Lw$ stretches the vehicle without changing its
      cross-section shape at the base; aerodynamic coefficients (which
      are non-dimensional against a fixed $\SI{1}{m^{2}}$ reference
      area) rise.
\item Increasing $W$ or decreasing $\Ls$ sharpens the spanwise Mach
      gradient and the outboard quartic shock, increasing the curvature
      near the tips.
\item Narrowing the gap $y_{6}-y_{5}$ flattens the upper surface and
      shrinks~$\Vol$.
\item Non-zero $\delta_{5},\delta_{6}$ skew the upper trailing edge.
\end{itemize}

% ==============================================================
\section{Validation Against the Paper}
% ==============================================================

The package's \texttt{validate.py} reproduces Table~4 of the paper for
the baseline (\mbox{$\Ma_{c}=6$}, $\Ma_{t}=13$) configuration and
compares $C_{L}, C_{D}, L/D, C_{m,z}, X_{\mathrm{cp}}$ across the
four-point Table~2 trajectory.

\begin{center}
\begin{tabular}{@{}lccc@{}}
\toprule
Metric & Computed & Paper & Deviation \\
\midrule
$\Vol$   [\si{m^{3}}]  & $3.42$   & $3.02$   & $+13.2\%$ \\
$\Swet$  [\si{m^{2}}]  & $23.25$  & $26.23$  & $-11.4\%$ \\
$\Sp$    [\si{m^{2}}]  & $9.71$   & $9.59$   & $+1.3\%$  \\
$\Sb$    [\si{m^{2}}]  & $1.58$   & $1.42$   & $+11.1\%$ \\
$\eta$                 & $0.0976$ & $0.0797$ & $+22.5\%$ \\
\bottomrule
\end{tabular}
\end{center}

The deviations arise from the simplified \emph{vertical osculating
plane} approximation adopted in the construction: the prompt's Step~4
reduces each plane's streamline to the straight line
Eq.~\eqref{eq:yTE}, which ignores the tilt of the true osculating
plane toward the local curvature centre for $|z|>\Ls$. Reaching the
paper's $\pm 3\%$ tolerance would require replacing
Eqs.~\eqref{eq:xLE}--\eqref{eq:yTE} with full 3D curvature-aware
osculating geometry (Rodi~2011); that enhancement is out of scope for
the present reproduction.

The aerodynamic trends are monotonic in~$\Ma$ and $X_{\mathrm{cp}}$ is
stable at $\sim 0.63$ (paper $\sim 0.678$), which is the principal
correctness indicator the paper highlights in Fig.~12(e).

% ==============================================================
\section{Command-Line Usage}
% ==============================================================

\begin{lstlisting}
# Validate against Table 4 + Fig. 12
python -m liu2019 --validate

# Generate geometry only and write an STL
python -m liu2019 --generate --stl out.stl --n-z 200 --n-x 100

# Skip the aero sweep during validation
python -m liu2019 --validate --no-aero
\end{lstlisting}

From Python:
\begin{lstlisting}
from liu2019 import build_liu2019_waverider, PAPER_PARAMS
wr = build_liu2019_waverider(PAPER_PARAMS, n_z=200, n_x=100)
print(wr.summary())
wr.export_stl("liu2019.stl")
\end{lstlisting}

% ==============================================================
\section{References}
% ==============================================================

\begin{enumerate}
\item J.~Liu, Z.~Liu, X.~Wen, F.~Ding, ``Novel osculating flowfield
      methodology for wide-speed range waverider vehicles across
      variable Mach number,'' \textit{Acta Astronautica} \textbf{162},
      160--167, 2019. DOI: \texttt{10.1016/j.actaastro.2019.05.056}.
\item H.~Sobieczky, F.~Dougherty, K.~Jones, ``Hypersonic waverider
      design from given shock waves,'' First International Waverider
      Symposium, University of Maryland, 1990.
\item V.~Rodi, ``The osculating flowfield method of waverider
      geometry generation,'' AIAA Paper 2011-1188, 2011.
\item J.~D.~Anderson, \textit{Modern Compressible Flow with Historical
      Perspective}, 3rd ed., McGraw-Hill, 2003 (Chs.~9--10).
\item P.~L.~Rudd and M.~J.~Lewis, ``Comparison of shock calculation
      methods,'' \textit{Journal of Aircraft} \textbf{35}(3),
      520--522, 1998.
\end{enumerate}

\end{document}
